\chapter{Considerações Finais}

Este trabalho teve como objetivo propor e desenvolver um protótipo inicial de
Avaliador de Serviços Públicos digitais, capaz de registrar sinais de dificuldade
na navegação, solicitar feedback contextual de baixo esforço e apresentar
indicadores agregados para apoiar a melhoria dos serviços avaliados. A proposta
partiu de um problema prático: avaliações exibidas apenas ao final da jornada nem
sempre capturam o momento em que o usuário se perde, demora ou abandona a busca
por um serviço.

A revisão teórica indicou que a avaliação de serviços públicos precisa combinar
satisfação, qualidade percebida, efetividade, usabilidade e orientação para
melhoria. Também reforçou que uma nota geral, isoladamente, explica pouco sobre o
problema vivido pelo cidadão. Por isso, o protótipo combina dois tipos de dado:
sinais comportamentais, como cliques e tempo de permanência, e feedback
declarado, como estrelas, atributos do problema, NPS e comentário opcional.

A Lean Inception contribuiu para transformar a ideia inicial em um MVP mais
delimitado. As dinâmicas de escopo, objetivos, personas, jornadas,
brainstorming, revisão técnica, sequenciamento e Canvas MVP permitiram separar o
que deveria entrar na primeira versão do que deveria permanecer como evolução.
Com isso, o trabalho evitou ampliar o escopo para ouvidoria, chatbot,
identificação individual, autenticação avançada ou integração real com portais
governamentais.

O protótipo desenvolvido demonstrou a viabilidade técnica do fluxo proposto. A
solução local registra eventos de navegação, aplica regras explícitas de
detecção, exibe popup de redirecionamento para serviços prioritários, abre uma
página de feedback e apresenta dados administrativos agregados. A contribuição do
artefato está em mostrar como o momento da dificuldade pode ser associado ao
serviço acessado, ao tipo de sinal observado e ao motivo declarado pelo usuário.

\section{Resultados em Relação aos Objetivos}

O primeiro objetivo específico foi levantar fundamentos legais, técnicos e
teóricos sobre avaliação de serviços públicos, qualidade percebida, usabilidade,
governo digital e Lean Inception. Esse objetivo foi atendido no referencial
teórico, que reuniu legislação brasileira, modelos de avaliação de serviços,
BRASP, API de avaliação, modelo unificado, experiência do usuário, NPS, métricas
de navegação, LGPD e fundamentos de Design Science Research.

O segundo objetivo específico foi definir, por meio de Lean Inception, o fluxo de
detecção de dificuldade, sugestão de redirecionamento e coleta de feedback. Esse
resultado aparece na definição do MVP, nas personas, nas jornadas, no
sequenciamento de funcionalidades e no fluxo conceitual apresentado nos
apêndices. A Lean Inception ajudou a justificar por que o produto deveria ser um
mecanismo contextual de apoio à melhoria, e não uma ouvidoria, um chatbot ou um
sistema completo de atendimento.

O terceiro objetivo específico foi construir um protótipo funcional com API,
páginas de serviço público, feedback contextual, formulário estático comparativo
e visualização administrativa. Esse objetivo foi atendido pela implementação do
servidor local, dos endpoints, das páginas navegáveis, do popup de
redirecionamento, da tela de feedback com estrelas, atributos, NPS e comentário,
e do painel administrativo com alertas e indicadores agregados.

O quarto objetivo específico foi estruturar os dados e indicadores necessários
para apoiar a análise de cliques, tempo de navegação, taxa de resposta,
avaliações, NPS e pontos recorrentes de melhoria. O protótipo já organiza esses
dados em eventos de clique, eventos de tempo, alertas administrativos e
feedbacks estruturados. Também disponibiliza endpoints de consulta para o painel.
A comparação entre avaliação estática e feedback contextual, entretanto, ainda
não possui resultado estatístico, pois depende de coleta com participantes em
cenário controlado. Nesta etapa, o resultado concreto foi deixar a estrutura
técnica preparada para medir essa comparação em uma validação posterior.

\section{Limitações}

A principal limitação é a ausência de validação com uma amostra real de usuários
finais. O protótipo permite executar o fluxo completo, mas ainda não mede, com
base empírica suficiente, se o feedback contextual aumenta a taxa de resposta ou
melhora a qualidade dos dados quando comparado a um formulário estático.

O desenvolvimento também evidencia uma limitação organizacional. A tecnologia,
isoladamente, não garante melhoria de serviços públicos. Para que o Avaliador
gere impacto real, o órgão precisa tratar feedback negativo como insumo de
aprendizado, definir responsáveis pela análise dos indicadores e transformar os
alertas em decisões de melhoria. Sem essa rotina institucional, o painel poderia
se tornar apenas um repositório de sinais de dificuldade e avaliações, sem
efeito concreto na experiência do cidadão.

Outra limitação está nos limiares de clique e tempo. Os valores usados no
backend são operacionais e servem para demonstrar a lógica de detecção:
dificuldade por clique acima do limite configurado e dificuldade por tempo acima
de 1,5 vez o tempo médio definido para o serviço. Em uma implantação real, esses
limiares deveriam ser calibrados com dados históricos, testes de usabilidade e
análise por tipo de serviço.

Há ainda um desafio ético e de privacidade. Mesmo sem coletar CPF, nome, e-mail
ou login, padrões de navegação, tempo de permanência, sequência de ações e tipo
de serviço acessado podem gerar inferências sobre comportamentos de grupos de
usuários. Por isso, qualquer implantação real deve aplicar rigorosamente os
princípios da LGPD, especialmente minimização, finalidade, transparência,
controle de acesso, retenção limitada e anonimização efetiva \cite{brasil2018lgpd}.

Também não houve integração com portais governamentais reais. Essa decisão
manteve o escopo compatível com o TCC e evitou dependências de autorização
institucional, segurança, governança e privacidade. Assim, os resultados devem
ser interpretados como demonstração de viabilidade conceitual e técnica, não como
comprovação de impacto em produção.

\section{Trabalhos Futuros}

Os trabalhos futuros devem avançar em três frentes. A primeira é a validação com
usuários. Para isso, recomenda-se organizar um estudo controlado com dois fluxos:
no fluxo A, o usuário navega e recebe apenas um formulário estático ao final; no
fluxo B, o sistema monitora sinais de dificuldade e solicita feedback contextual.
Essa comparação deve medir taxa de resposta, tempo até a resposta, quantidade de
atributos selecionados, qualidade dos comentários, taxa de abandono e clareza dos
dados gerados para o gestor. Só depois dessa coleta será possível afirmar, com
base empírica, se o feedback contextual supera o formulário estático no cenário
avaliado.

A segunda frente é a calibração técnica. Os limiares atuais de clique e tempo
foram definidos para demonstrar o funcionamento do protótipo. Em uma evolução, o
sistema deve permitir configuração por serviço, por página e por tipo de tarefa.
Também deve registrar versão das regras, evitar alertas duplicados, separar
eventos por organização, criar filtros por período e serviço, e apresentar
métricas comparáveis no painel. Essa etapa aproxima o protótipo de uma solução
mais robusta de engenharia de software, com regras auditáveis e manutenção mais
segura.

A terceira frente é a preparação para piloto autorizado. Uma implantação real
exigiria aviso de coleta, definição de base legal e finalidade, política de
retenção, controle de acesso ao painel, chave de API, separação por órgão,
monitoramento de disponibilidade e documentação de integração do script cliente.
Também seria necessário validar se os indicadores gerados realmente apoiam a
decisão administrativa, por exemplo, priorização de páginas confusas, revisão de
linguagem, redução de etapas e reorganização de menus.

\section{Conclusão}

Conclui-se que o Avaliador de Serviços Públicos proposto é uma alternativa
viável para explorar feedback contextual em serviços digitais. A solução não
substitui instrumentos oficiais de avaliação, mas pode complementá-los ao
capturar sinais de dificuldade durante a navegação e relacioná-los a avaliações
simples feitas pelo usuário. Essa combinação responde ao problema inicial do
trabalho: avaliações genéricas e tardias podem indicar insatisfação, mas nem
sempre mostram onde a jornada falhou ou qual característica precisa ser
melhorada.

O artefato desenvolvido mostra que é tecnicamente possível registrar eventos de
navegação em uma sessão anônima, aplicar regras de detecção, sugerir
redirecionamento em serviços prioritários, solicitar feedback de baixo esforço e
organizar os dados em indicadores administrativos. Mesmo sem validação
estatística ampla nesta etapa, o protótipo deixa pronta uma base mensurável para
testar a hipótese em etapa posterior.

O percurso do trabalho também mostra que a solução não nasce apenas de uma
decisão técnica de implementação. Ela foi delimitada por fundamentos de governo
digital, qualidade percebida, usabilidade, privacidade e Lean Inception, e por
artefatos de produto que explicitaram personas, jornadas, escopo, features,
sequenciamento e critérios de verificação. Isso torna a proposta mais defensável
como trabalho de Engenharia de Software, pois conecta problema, método, protótipo
e indicadores.

A principal contribuição do trabalho está na integração entre fundamentos de
avaliação de serviços públicos, Lean Inception, design centrado no usuário,
prototipação de software e uma API demonstrável para análise da experiência
digital. A proposta torna mais concreta a ideia de que a melhoria de serviços
públicos digitais pode ser apoiada por dados de contexto, desde que esses dados
sejam coletados com finalidade clara, baixo esforço para o cidadão e respeito à
privacidade.
